\SourceLeaf{072}
\(Y,Z\) etc. hinc ad valorem evanescentem reductum iri. Simili ratione ex aequatione (11) colligitur, \(K\) esse valorem functionis \(u\) ex eadem substitutione emergentem.

Applicationem correctionum \(-\varepsilon,-\varepsilon',-\varepsilon''\) etc. ad observationes, vocabimus \textit{observationum compensationem}, manifestoque deducti sumus ad conclusionem gravissimam, puta, observationes eo quem docuimus modo compensatas omnibus aequationibus conditionalibus exacte satisfacere, atque cuilibet quantitati ab observationibus quomodocunque pendenti eum ipsum valorem conciliare, qui ex observationum non mutatarum combinatione maxime idonea emergeret. Quum itaque impossibile sit, errores ipsos \(e,e',e''\) etc. ex aequationibus conditionalibus eruere, quippe quarum multitudo haud sufficit, saltem \textit{errores maxime plausibiles} nacti sumus, qua denominatione quantitates \(\varepsilon,\varepsilon',\varepsilon''\) etc. designare licebit.

\Article{10}
Quum multitudo observationum maior esse supponatur multitudine aequationum conditionalium, praeter systema correctionum maxime plausibilium \(-\varepsilon,-\varepsilon',-\varepsilon''\) etc. infinite multa alia inveniri possunt, quae aequationibus conditionalibus satisfaciant, operaeque pretium est indagare, quomodo haec ad illud se habeant. Constituant itaque \(-E,-E',-E''\) etc. tale systema a maxime plausibili diversum, habebimusque
\[
\begin{aligned}
aE+a'E'+a''E''+\text{etc.}&=\mathfrak A\\
bE+b'E'+b''E''+\text{etc.}&=\mathfrak B\\
cE+c'E'+c''E''+\text{etc.}&=\mathfrak C
\end{aligned}
\]
etc. Multiplicando has aequationes resp. per \(A,B,C\) etc. et addendo, obtinemus adiumento aequationum (10)
\[
p\varepsilon E+p'\varepsilon'E'+p''\varepsilon''E''+\text{etc.}=A\mathfrak A+B\mathfrak B+C\mathfrak C+\text{etc.}
\]
Prorsus vero simili modo aequationes (13) suppeditant
\[
p\varepsilon\varepsilon+p'\varepsilon'\varepsilon'+p''\varepsilon''\varepsilon''+\text{etc.}=A\mathfrak A+B\mathfrak B+C\mathfrak C+\text{etc.}\quad(14)
\]
E combinatione harum duarum aequationum facile deducitur
\[
\begin{aligned}
&pEE+p'E'E'+p''E''E''+\text{etc.}\\
&=p\varepsilon\varepsilon+p'\varepsilon'\varepsilon'+p''\varepsilon''\varepsilon''+\text{etc.}
\\&\qquad+p(E-\varepsilon)^2+p'(E'-\varepsilon')^2+p''(E''-\varepsilon'')^2+\text{etc.}
\end{aligned}
\]
\SourceLeaf{073}
Aggregatum \(pEE+p'E'E'+p''E''E''+\text{etc.}\) itaque necessario \textit{maius} erit aggregato \(p\varepsilon\varepsilon+p'\varepsilon'\varepsilon'+p''\varepsilon''\varepsilon''+\text{etc.}\), quod enuntiari potest tamquam

\textsc{Theorema.} \textit{Aggregatum quadratorum correctionum, per quas observationes cum aequationibus conditionalibus conciliare licet, per pondera observationum resp. multiplicatorum, fit minimum, si correctiones maxime plausibiles adoptantur.}

Hoc est ipsum principium quadratorum minimorum, ex quo etiam aequationes (12), (10) facile immediate derivari possunt. Ceterum pro hoc aggregato minimo, quod in sequentibus per \(S\) denotabimus, aequatio (14) nobis suppeditat expressionem \(\mathfrak A A+B\mathfrak B+C\mathfrak C+\text{etc.}\)

\Article{11}
Determinatio errorum maxime plausibilium, quum a coëfficientibus \(l,l',l''\) etc. independens sit, manifesto praeparationem commodissimam sistit, ad quemvis usum, in quem observationes vertere placuerit. Praeterea perspicuum est, ad illud negotium haud opus esse eliminatione \textit{indefinita} seu cognitione coëfficientium \([\alpha\alpha],[\alpha\beta]\) etc., nihilque aliud requiri, nisi ut quantitates auxiliares \(A,B,C\) etc., quas in sequentibus \textit{correlata} aequationum conditionalium \(X=0,\ Y=0,\ Z=0\) etc. vocabimus, ex aequationibus (12) per eliminationem definitam eliciantur atque in formulis (10) substituantur.

Quamquam vero haec methodus nihil desiderandum linquat, quoties quantitatum ab observationibus pendentium valores maxime plausibiles tantummodo requiruntur, tamen res secus se habere videtur, quoties insuper pondus alicuius determinationis in votis est, quum ad hunc finem, prout hac vel illa quatuor expressionum supra traditarum uti placuerit, cognitio quantitatum \(L,L',L''\) etc., vel saltem cognitio harum \(x^0,y^0,z^0\) etc. necessaria videatur. Hac ratione utile erit, negotium eliminationis accuratius perscrutari, unde via facilior ad pondera quoque invenienda se nobis aperiet.

\Article{12}
Nexus quantitatum in hac disquisitione occurrentium haud parum illustratur per introductionem functionis indefinitae secundi ordinis
\[
\begin{aligned}&[aa]xx+2[ab]xy+2[ac]xz+\text{etc.}\\&\qquad+[bb]yy+2[bc]yz+\text{etc.}+[cc]zz+\text{etc.}\end{aligned}
\]
quam per \(T\) denotabimus. Primo statim obvium est, hanc functionem fieri
